\chapter{Glossary of Terms}
\label{ch:glossary}

\begin{description}[style=nextline, leftmargin=2em, font=\bfseries\color{navyblue}]

\item[1-9-90 Pattern]
The predictable adoption distribution identified by BCG in AI deployment: 1 per cent are natural innovators who adopt independently, 9 per cent become champions with organisational support (the judgment broker candidates), and 90 per cent follow once social proof exists. Strategic investment in the 9 per cent yields the highest return on adoption.

\item[Agent Colony]
An organisational model (associated with Every.to) in which every employee is given a persistent, personalised AI agent. Coordination emerges from cultural norms and public visibility rather than architectural design. Proven at approximately 20-person scale; subject to the ant death spiral at larger scales.

\item[Agentic Mesh]
See \emph{Dynamic Agentic Mesh}.

\item[Ant Death Spiral]
A failure mode in multi-agent systems where agents respond to each other's outputs in an uncontrolled loop, consuming resources (tokens, compute, time) without producing useful outcomes. Named by analogy with army ant death spirals in entomology. Documented in Every's agent colony experiments at 20-person scale.

\item[Augmentation Ratio (AR)]
A proposed KPI measuring the proportion of decision-making volume successfully handled by agents without requiring human escalation. Target: $>65\%$. Too low indicates the mesh is not yet effective; too high indicates potential vigilance degradation.

\item[BYOAI]
Bring Your Own AI --- the practice of employees using personal AI tools and accounts for work purposes, bypassing corporate-approved systems. Microsoft/LinkedIn 2024 data shows 78 per cent of knowledge workers engage in BYOAI.

\item[Centaur]
A human-AI collaboration mode in which human and AI work with clear division of labour: the human performs tasks where human judgment is strongest, the AI handles tasks inside the frontier, with distinct handoff points. One of the two high-performing modes identified by Randazzo et al.\ (2024).

\item[Compound Engineering]
A principle articulated by Kieran Klaassen (Every.to, 2025): each unit of engineering work should make subsequent units easier, not harder. The 50/50 Rule: spend half your time building features, half improving the system that builds features. The organisational analogue is the compounding organisation.

\item[Compounding Organisation]
An organisation in which each unit of organisational learning makes the next one easier. The flywheel proceeds: workforce discovers a shortcut, agents detect and codify it, the judgment broker validates it, the pattern propagates across the mesh, and the organisation's capability compounds.

\item[Coordination Tax]
The cost --- in headcount, latency, distortion, and overhead --- of maintaining hierarchical layers whose primary function is routing information between organisational levels. This report argues that AI collapses the cost of the coordination tax to near-zero.

\item[Cyborg]
A human-AI collaboration mode characterised by continuous interleaving: human and AI work on the same analysis, passing work back and forth across the frontier boundary without clean separation of responsibilities. One of the two high-performing modes identified by Randazzo et al.\ (2024).

\item[DAG (Directed Acyclic Graph)]
A graph structure in which edges have direction and no cycles exist. In the Dynamic Agentic Mesh, DAGs represent formalised workflows: sequences of tasks with dependencies, ownership, and governance constraints. Shadow workflows are codified as DAGs when promoted into the mesh.

\item[Declarative Governance]
A governance model in which policies are expressed as code and enforced automatically at every system transition, rather than being applied through periodic human review. Analogous to declarative infrastructure (Kubernetes, Terraform) in software engineering.

\item[Discovery Engine]
The bottom-up layer of the Dynamic Agentic Mesh that captures organisational insights. Passive agent monitoring surfaces shadow workflows; natural-language Insight Portals allow anyone to propose improvements; tacit knowledge is codified into agentic instructions.

\item[Dynamic Agentic Mesh]
The organisational architecture proposed in this report. Three layers (Discovery Engine, Orchestration Layer, Declarative Governance) retain human judgment at intersection points while agentic orchestration handles information routing. Middle managers evolve into judgment brokers.

\item[HITL (Human-in-the-Loop)]
The practice of maintaining human oversight or decision-making authority within an automated system. This report argues that na\"{\i}ve HITL implementations degrade as AI quality improves (the vigilance problem) and that HITL must be architecturally supported through deliberate friction and active engagement.

\item[HITL Precision (HP)]
A proposed KPI measuring the proportion of escalated cases that genuinely required human intervention. Target: $>90\%$. Low precision indicates alert fatigue; high precision indicates well-calibrated mesh boundaries.

\item[Information Routing Protocol]
This report's characterisation of hierarchy: a system designed to filter, aggregate, and relay information across organisational scales, constrained by the human cognitive limit of managing three to eight direct reports.

\item[Insight Ingestion Loop]
The five-stage process by which the Dynamic Agentic Mesh converts individual discoveries into organisational capability: Discovery, Codification, Validation, Integration, Amplification.

\item[Intelligence Layer]
The organisational model proposed by Jack Dorsey and Roelof Botha (Block, March 2026) in which middle management is eliminated and coordination is handled by AI world models (Company World Model and Customer World Model) composed by an Intelligence Layer.

\item[Jagged Frontier]
The unpredictable boundary between tasks that AI performs well and tasks that AI performs poorly. Coined by Dell'Acqua et al.\ (2023/2026) based on a preregistered experiment with 758 BCG consultants. The frontier is ``jagged'' because adjacent, similar-looking tasks can fall on opposite sides.

\item[Judgment Broker]
The evolved middle management role proposed in this report. The judgment broker operates at the jagged frontier, providing strategic direction, ethical adjudication, and relational intelligence that agentic systems cannot perform. Distinguished from the traditional information-routing role by its focus on edge cases, cross-functional coherence, and frontier probing.

\item[Mesh Velocity (MV)]
A proposed KPI measuring the elapsed time from insight discovery to codified, sanctioned workflow in the mesh. Target: $<48$ hours. The primary indicator of organisational learning speed.

\item[Ontological Provenance]
The property of an agentic system in which every decision traces back through a formal knowledge graph (OWL~2 ontology), providing auditable reasoning chains from output to input, policy, and human decision.

\item[OWL~2]
Web Ontology Language 2 --- a W3C standard for formal knowledge representation. In VisionClaw, OWL~2 provides shared semantics for agent skill routing through ontological subsumption (understanding that ``risk assessment'' is a sub-task of ``governance review'').

\item[Secret Cyborg]
A term coined by Ethan Mollick (2023) for workers who use AI tools without disclosing this use to their employer. The secrecy is driven by rational calculation: disclosure risks headcount reduction if the worker demonstrates they have automated significant portions of their role.

\item[Self-Automator]
A human-AI collaboration mode in which the human delegates wholesale to AI with minimal oversight. The lowest-performing mode in the Randazzo et al.\ (2024) study. Block's organisational model is characterised in this report as creating Self-Automator organisations.

\item[Shadow AI]
AI tools and applications used by employees without organisational awareness, approval, or governance. Shadow AI usage increased 156 per cent between 2023 and 2025; the average enterprise hosts approximately 1,200 unauthorised AI applications.

\item[Shadow Workflow]
An informal work process created by an employee using AI tools, which bypasses or replaces the officially documented workflow without management awareness. Shadow workflows represent both a risk (ungoverned, undocumented) and an asset (raw material for the compounding loop).

\item[Trust Variance (TV)]
A proposed KPI measuring the standard deviation of agent decision quality over a rolling 30-day window. Target: $<0.12\sigma$. The early warning system for vigilance degradation --- catching drift before it compounds.

\item[Vigilance Problem]
The empirically demonstrated tendency of human oversight quality to degrade as AI quality improves. Identified by Dell'Acqua (2023) in an experiment with HR recruiters: those given high-quality AI made worse decisions than those given low-quality AI because they stopped scrutinising outputs.

\item[VisionClaw]
The operational agentic mesh system built by DreamLab AI Consulting. Architecture: OWL~2 ontology engine, Claude-Flow orchestration (83 skills), WebXR 3D visualisation (CUDA-accelerated), Nostr agent identity, RuVector memory (1.17M+ vector embeddings). Proven at 15-person scale.

\end{description}
